

Our work focuses on enhancing the semantic accuracy of LLMs through constrained decoding. 
Prior research has explored two primary strategies to improve LLM accuracy in generating structured formal languages:
Fine-tuning or prompt engineering \cite{bassamzadeh2024comparativestudydslcode, weyssow2024exploringparameterefficientfinetuningtechniques}, which typically requires significant data, computational resources, and time, often without formal guarantees of success.
However, fine-tuning and prompt engineering approaches are complementary to the constrained decoding approach we adopt, and improvements from those techniques could enhance the overall quality of LLM output.


Context-free-grammar generation techniques such as \tgcd{}~\cite{geng2023grammar}, \outlines{}~\cite{willard2023efficient}, \domino{}~\cite{beurerkellner2024guiding}, \syncode{}~\cite{ugare2024syncodellmgenerationgrammar} and \aici{}~\cite{Moskal2024} constrain LLM output according to grammar rules. 
However, in contrast to \Tool{}, these tools cannot apply semantic constraints to the generation process.
% \sasa{check the rewrite from but do not apply semanitc to this.}
Other recent constrained-generation methods utilize language servers (designed for communication between IDEs and language-specific tools) to enforce some semantic constraints during decoding \cite{agrawal2023guiding, Wei_2023}. 
However, these techniques lack guarantees for syntactic accuracy and depend on the availability and performance of language servers.
% \sasa{What is lang server tool? seems a convoluted way to say "code completion tools"}

% \\synchromesh{}, \guidance{} and \lmql{}
% \lmql{}~\cite{10.1145/3591300} is a query language designed for structured LLM generation, allowing users to write queries with "holes" in text, where constraints can be applied. 
% % \lmql{} fills these holes using rejection sampling based on the constraints. 
% These constraints are limited to regular expressions or predefined types such as int and float and do not extend to context-free grammars.
% In contrast, \Tool{} operates on a predefined overarching context-free grammar and offers fine-grained control over the generation process to the user. 
% Users can apply rejection sampling to specific parts of the grammar by moving forward or backward through the output. 
% Additionally, \Tool{} allows users to adjust generation parameters, providing flexibility during generation dynamically.

\guidance{}~\cite{guidance} supports context-free languages but requires users to compose grammars through supported operations. 
\guidance{}’s \str{stop\_at} function, which halts generation at a specified regular expression, has similarities to the \Tool{}’s \str{forward} function. 
However, while \str{stop\_at} works with regular expressions, \str{forward} operates based on symbols from \Tool{}'s overarching grammar. 
Unlike \Tool{}, \guidance{} does not support backtracking, and the only way to impose constraints is through regular expressions on generated "holes," similar to \lmql{}. 
Moreover, \Tool{} uses any LR grammar in the standard Lark EBNF format, making it easier to plug in large grammars like SQL, which is not straightforward with \guidance{}.
Both \lmql{} and \guidance{} provide additional features, such as the ability to insert strings during generation and support for function calls, which are outside the scope of this paper.

\synchromesh{}~\cite{poesia2022synchromesh} uses constrained semantic decoding (CSD) to enforce semantic constraints through predictive masking and rejection sampling at the token level. 
It checks if the model's first token choice adheres to the semantic constraints, and if not, uses predictive masking to resample. 
It is designed for use with OpenAI's GPT-3 and Codex and relies on API access without direct control over the underlying language models. 
Similarly, \picard{}~\cite{scholak-etal-2021-picard} is a grammar-guided generation tool that's developed for SQL generation with additional constraints on valid table and column names. 
The approach used in \synchromesh{} and \picard{} for SQL can be easily implemented with \Tool{} with few lines of code, as shown in our case study. 
In contrast to both \synchromesh{} and \picard{}, the goal of \Tool{} is to develop an efficient and intuitive tool that allows users to write programs to define grammar-level semantic constraints through its forward and backward operations that can work with any user-provided grammar and not specific to improving SQL generation. \add{An unofficial implementation of Synchromesh exists; in practice, this system encountered errors when running with complex Lark grammars}. Furthermore, \picard{} works only with T5 architecture, and thus it is not possible to make an empirical comparison to \Tool{}.

\add{
\Tool{} also serves as the primary building block in recent works like CRANE \cite{banerjee2025crane}, which combines syntactic and semantic correctness of constrained decoding with unconstrained LLM reasoning steps to improve LLM performance on challenging symbolic reasoning benchmarks such as GSM-symbolic and FOLIO.
}

% \sasa{Is there some other principled thing to shoot picard down? E.g. it works only for sql, single architecure; vs also stating Itergen is general in both dimensions (any cfg language and supports various archs.)}
